\documentclass[11pt]{article}

\usepackage[margin=1in]{geometry}
\usepackage{array}
\usepackage{booktabs}
\usepackage{enumitem}
\usepackage{xcolor}
\usepackage{colortbl}
\usepackage{tikz}
\usetikzlibrary{arrows.meta, positioning, shapes.geometric, fit, backgrounds, matrix}
\usepackage{pdflscape}
\usepackage{hyperref}
\usepackage{titlesec}

\definecolor{mathcol}{HTML}{1F77B4}
\definecolor{cscol}{HTML}{2CA02C}
\definecolor{phycol}{HTML}{D62728}
\definecolor{eececol}{HTML}{9467BD}
\definecolor{ewcol}{HTML}{FF7F0E}
\definecolor{econcol}{HTML}{8C564B}
\definecolor{phaseband}{HTML}{EFEFEF}

\titleformat{\section}{\Large\bfseries}{\thesection}{0.6em}{}[\vspace{-0.4em}\rule{\linewidth}{0.4pt}]
\titleformat{\subsection}{\large\bfseries}{\thesubsection}{0.6em}{}

\newcommand{\track}[2]{\textcolor{#1}{\textbf{#2}}}
\newcommand{\book}[1]{\textit{#1}}

\title{\vspace{-2em}Self-Taught Engineering Curriculum\\\large A Progressive, Multi-Track Reading Plan}
\author{}
\date{}

\begin{document}
\maketitle
\vspace{-3em}

\section*{Overview}

This curriculum weaves six subject tracks into a progressive sequence of
\emph{phases}. Within a phase, tracks run in parallel --- read from multiple
books at once. Between phases, prerequisite dependencies force an order
(e.g. differential equations and Fourier analysis must precede DSP;
digital systems and C must precede computer organization).

\smallskip
\noindent\textbf{Tracks}:
\track{mathcol}{Mathematics},
\track{cscol}{Computer Science},
\track{phycol}{Physics},
\track{eececol}{Electrical \& Computer Engineering},
\track{ewcol}{Electronic Warfare},
\track{econcol}{Economics}.

\section{The Six Phases}

\renewcommand{\arraystretch}{1.2}
\begin{center}
{\small
\begin{tabular}{>{\bfseries\raggedright\arraybackslash}p{0.22\linewidth}
                >{\raggedright\arraybackslash}p{0.40\linewidth}
                >{\raggedright\arraybackslash}p{0.32\linewidth}}
\toprule
\textbf{Phase} & \textbf{Theme} & \textbf{Unlocks} \\
\midrule
\rowcolor{phaseband} 0 -- Foundations & Numeracy, algebra, first programs. & Calculus, all computing tracks. \\
1 -- Core Tools   & Calculus I--III, C, classical mechanics. & Linear algebra, circuits, DSA. \\
\rowcolor{phaseband} 2 -- Structure & Linear algebra, discrete math, probability, DSA, digital logic, E\&M physics. & Diff.\ eq., circuits, OS, architecture. \\
3 -- Applied Core & ODEs, numerical methods, circuits fundamentals, OS, computer org. & Signals, control, embedded, microelectronics. \\
\rowcolor{phaseband} 4 -- Specialization & Signals \& systems, microelectronics, semiconductors, EM fields, embedded, VHDL, random processes. & DSP, comms, control, power, radar, EW. \\
5 -- Advanced Eng. & DSP, communications, feedback control, power systems. & Applied EW, radar, distributed systems. \\
\rowcolor{phaseband} 6 -- Applied Domain & Electronic warfare sequence, airborne radar, distributed systems. & --- \\
\bottomrule
\end{tabular}
}
\end{center}

\medskip
\noindent Engineering Economics sits outside the dependency graph --- read it
in any phase after basic algebra; it pairs well with Phase~3 as a
professional-practice topic.

\newpage
\section{Phase-by-Phase Reading List}

\subsection*{Phase 0 --- Foundations}
Goal: algebra fluency, first programs, habits of study.
\begin{itemize}[leftmargin=1.2em, itemsep=2pt]
  \item \track{mathcol}{Math} --- \book{Using and Understanding Mathematics:
        A Quantitative Reasoning Approach} (Math 102) \\
        \book{Algebra: Form and Function} (Math 103) \\
        \book{Pre-Calculus, 2e} (Math 109/110)
  \item \track{cscol}{CS} --- \book{Python for Everyone} (CMPS 150) \\
        \book{Introduction to Programming Using Python} (CMPS 150)
\end{itemize}

\subsection*{Phase 1 --- Core Tools}
Goal: the three pillars --- calculus, systems programming, classical physics.
\begin{itemize}[leftmargin=1.2em, itemsep=2pt]
  \item \track{mathcol}{Math} --- \book{Technical Mathematics with Calculus}
        (Math 210) \emph{or} \book{Applied Calculus} (Math 250) as a gentler
        entry; use \book{Calculus: Early Transcendentals, 11e} (Math 270/301/302)
        as the authoritative text through multivariable.
  \item \track{cscol}{CS} --- \book{The C Programming Language} (K\&R), then
        \book{C Programming: A Modern Approach}. (In parallel with Math.)
  \item \track{phycol}{Physics} --- \book{Fundamentals of Physics, 11e}
        (mechanics chapters); 10e Extended covers the same material if preferred.
\end{itemize}

\newpage
\subsection*{Phase 2 --- Structure}
Goal: the discrete, linear, and probabilistic scaffolding underneath
engineering and CS; finish classical physics; cross into hardware and
data structures.
\begin{itemize}[leftmargin=1.2em, itemsep=2pt]
  \item \track{mathcol}{Math} --- \book{Elementary Linear Algebra, 11e}
        (Math 362) \\
        \book{Probability and Statistics with R} (Math 325) \\
        \book{Discrete Structures, Logic, and Computability, 4e} (CMPS 341)
  \item \track{cscol}{CS} --- \book{Data Structures and Algorithm Analysis
        in C++} \\
        \book{Introduction to Java Programming and Data Structures}
        (CMPS 260/261) for OOP exposure
  \item \track{phycol}{Physics} --- finish \book{Fundamentals of Physics}
        (E\&M, waves, optics); \book{Statics and Mechanics of Materials, 5e}
  \item \track{eececol}{EECE} --- \book{Digital Systems, 12e} (EECE 240) ---
        only algebra and Boolean logic are needed, so start it here.
\end{itemize}

\subsection*{Phase 3 --- Applied Core}
Goal: the mathematics of change (ODEs) and the engineering it enables.
Econ fits naturally here.
\begin{itemize}[leftmargin=1.2em, itemsep=2pt]
  \item \track{mathcol}{Math} --- \book{Elementary Differential Equations}
        (Math 350) \\
        \book{A Transition to Advanced Mathematics} (Math 360) --- optional,
        for proof-based math and rigor before later theory.
  \item \track{eececol}{EECE} --- \book{Fundamentals of Electric Circuits, 6e}
        (EECE 355/353) \\
        \book{Applied Numerical Methods with MATLAB for Engineers and
        Scientists, 3e} (EECE 260)
  \item \track{cscol}{CS} --- \book{Operating System Concepts} \\
        \book{Computer Organization and Design} (the CMPS edition) ---
        needs C and digital logic, which are now in hand.
  \item \track{econcol}{Econ} --- \book{Fundamentals of Engineering Economics}
        (read in parallel; no hard prerequisites beyond algebra).
\end{itemize}

\subsection*{Phase 4 --- Specialization}
Goal: the EECE body of knowledge --- devices, fields, signals, embedded
hardware. Several books are parallel; \book{Signals and Systems} is the
gateway to Phase~5.
\begin{itemize}[leftmargin=1.2em, itemsep=2pt]
  \item \track{eececol}{Microelectronics} --- \book{Microelectronic
        Circuits, 7e} (EECE 353), with \book{Laboratory Explorations for
        Microelectronic Circuits} and \book{The Art of Electronics, 3e}
        as a practitioner companion.
  \item \track{eececol}{Devices} --- \book{Fundamentals of Semiconductor
        Devices, 2e} (EECE 335).
  \item \track{eececol}{Fields} --- \book{Fundamentals of Applied
        Electromagnetics} (EECE 344).
  \item \track{eececol}{Embedded} --- \book{AVR Microcontroller and
        Embedded Systems Using Assembly and C} (EECE 340); pair with
        \book{The Designer's Guide to VHDL} for hardware description.
  \item \track{eececol}{Signals \& Randomness} ---
        \book{Signals and Systems, 3e} (EECE 444) is the keystone: read it
        \emph{after} ODEs and linear algebra, \emph{before} DSP, control,
        and communications. Alongside: \book{Probability, Random Variables
        and Stochastic Processes, 4e} (Papoulis, EECE 380 Book~1) with the
        friendlier \book{Probability and Stochastic Processes for EE/CE, 3e}
        (EECE 380 Book~2) as a warmup.
  \item \track{eececol}{Architecture} --- \book{Computer Organization and
        Design, 5e} (EECE 459) as a deeper revisit if desired.
\end{itemize}

\subsection*{Phase 5 --- Advanced Engineering}
Goal: the four classical EE capstone areas. Each depends on Signals \&
Systems; communications also requires probability; control requires ODEs.
\begin{itemize}[leftmargin=1.2em, itemsep=2pt]
  \item \track{eececol}{DSP} --- \book{The Scientist and Engineer's Guide
        to Digital Signal Processing, 2e} (EECE 430). Prereqs: Fourier
        analysis from Signals \& Systems; linear algebra; some complex
        analysis from calculus.
  \item \track{eececol}{Communications} --- \book{Digital and Analog
        Communication Systems, 4e} and \book{Modern Digital and Analog
        Communication Systems} (EECE 452). Prereqs: Signals \& Systems,
        probability/random processes.
  \item \track{eececol}{Control} --- \book{Feedback Control of Dynamic
        Systems, 7e} (EECE 461); the other two Franklin/Powell copies are
        earlier editions of the same text --- use one as the primary, the
        other for alternate problem sets. Prereqs: ODEs, Laplace
        transforms, linear algebra.
  \item \track{eececol}{Power} --- \book{Power Systems Analysis}
        (EECE 447/450). Prereqs: circuits, electromagnetics, some ODEs.
\end{itemize}

\subsection*{Phase 6 --- Applied Domain}
Goal: apply the stack to electronic warfare and distributed systems.
\begin{itemize}[leftmargin=1.2em, itemsep=2pt]
  \item \track{ewcol}{EW core (Adamy)} --- read strictly in order:
        \book{EW 101: A First Course in Electronic Warfare} $\to$
        \book{EW 102: A Second Course} $\to$
        \book{EW 103: Tactical Battlefield Communications EW} $\to$
        \book{EW 104: EW Against a New Generation of Threats} $\to$
        \book{EW 105: Space Electronic Warfare}.
  \item \track{ewcol}{Radar} --- \book{Introduction to Airborne Radar}
        after Signals \& Systems and some DSP; use the \book{Electronic
        Warfare and Radar Systems Handbook} as a reference throughout the
        entire EW track.
  \item \track{cscol}{Distributed} --- \book{Distributed Modeling and
        Analysis} as a late-stage CS read once OS and probability are
        comfortable.
\end{itemize}

\section{Critical Prerequisite Chains (``Do not skip'')}

\begin{enumerate}[leftmargin=1.4em, itemsep=2pt]
  \item \textbf{DSP chain}: Calculus $\to$ Linear Algebra $\to$
        Differential Equations $\to$ Signals \& Systems $\to$ DSP.
        Attempting DSP without Fourier analysis from Signals \& Systems
        is the single most common self-study failure point.
  \item \textbf{Communications chain}: Signals \& Systems \emph{and}
        Probability / Random Processes $\to$ Communications.
        Random processes are not optional; modulation, noise, and
        detection theory all depend on them.
  \item \textbf{Control chain}: Differential Equations \emph{and} Linear
        Algebra $\to$ Signals \& Systems (for Laplace and transfer
        functions) $\to$ Feedback Control.
  \item \textbf{Microelectronics chain}: Physics (E\&M) $\to$ Electric
        Circuits $\to$ Microelectronics; Semiconductor Devices reinforces
        the physical model and can be read in parallel with
        Microelectronics.
  \item \textbf{Computer systems chain}: C $\to$ Digital Systems $\to$
        Computer Organization $\to$ Operating Systems.
        Data Structures \& Algorithms runs in parallel with the middle
        of this chain.
  \item \textbf{Radar / EW chain}: Electromagnetics \emph{and} Signals \&
        Systems \emph{and} Communications $\to$ EW 101--105 and Airborne
        Radar. The Adamy series itself is sequential.
\end{enumerate}

\section{Suggested Weekly Rhythm}

Read from two or three tracks per week rather than finishing one book
before starting another. A sustainable pattern:

\begin{itemize}[leftmargin=1.2em, itemsep=2pt]
  \item One \track{mathcol}{math} text (active problem solving).
  \item One \track{eececol}{engineering} or \track{cscol}{CS} text
        currently in-phase (applied reading).
  \item One ``survey'' or reference read --- Art of Electronics, the EW
        Handbook, or Engineering Economics --- for breadth without the
        cognitive load of a new framework.
\end{itemize}

\noindent Review proofs and derivations by redoing them from a blank page
a week later; this is the single highest-leverage habit for retaining
mathematical material that later tracks depend on.

\end{document}
